\documentclass[10pt,letterpaper]{article}
\usepackage{cornell-notes}

\title{Clause 22.04.06: Calling A Function With A Tuple Of Arguments}
\author{}
\date{}
\setDocTitle{Clause 22.04.06: Calling A Function With A Tuple Of Arguments}
\setDocSubtitle{C++ 2024}
\setDocOwner{C++ 2024 Cornell Notes}
\setDocDate{\today}
\setCornellCollection{C++ 2024 Cornell Notes}
\setCornellUnitType{Clause}
\setCornellUnitNumber{22.04.06}
\setCornellUnitTitle{Calling A Function With A Tuple Of Arguments}
\hypersetup{pdftitle={Clause 22.04.06: Calling A Function With A Tuple Of Arguments},pdfsubject={Cornell notes},pdfauthor={}}

\begin{document}
\maketitle
\begin{CornellOverviewBox}{Overview}
\textbf{Classification:} Standard library - Normative\\[2pt]
\textbf{Learning objective:} Understand the rules, terminology, and semantic consequences associated with Calling a function with a tuple of arguments.
\end{CornellOverviewBox}\begin{CornellSummaryBox}{Summary}
{\sffamily\bfseries\color{Accent} Summary}\par\vspace{3pt}Section 22.4.6 focuses on Calling a function with a tuple of arguments. Apply the specified interface together with its mandates, effects, result conditions, exception behavior, and complexity constraints. When applying it, preserve the distinction between mandatory requirements, recommendations, permissions, and behavior deliberately left undefined, unspecified, or implementation-defined.
\end{CornellSummaryBox}

\end{document}
